\section{Introduction}
\label{sec:introduction}

Modern foundation models have shifted the performance frontier of AI inference
from arithmetic throughput to \emph{data coordination}.
When model weights, activations, and key--value (KV) caches exceed accelerator
memory, inference latency is governed less by compute and more by how efficiently
execution coordinates data movement across a hierarchical memory system spanning
device memory, host memory, and secondary storage~\cite{ma2026inference,gholami2024memorywall}.

This shift carries pressing economic and environmental stakes.
Once a model is trained, the recurring cost of \emph{serving} it dominates its
lifetime resource footprint.
Hardware analyses identify memory capacity, memory bandwidth, and interconnect
latency---not arithmetic throughput---as the binding constraints of the inference
phase~\cite{zhou2024survey}.
These constraints translate into macro-level exposure: aggregate capital
expenditure across major hyperscalers is on track to overtake operating cash
inflows by end of 2026~\cite{juniewicz2026hyperscaler}, and AI infrastructure
is absorbing available memory supply faster than it can be
produced~\cite{chey2026tpd}.
In this context, characterizing the \emph{intrinsic limits} of memory
orchestration is not merely a systems-optimization concern but a prerequisite
for the financial and energy viability of inference-driven AI.

A large body of systems work reduces memory pressure via CPU/NVMe offloading
and heterogeneous execution~\cite{flexgen2023,aminabadi2022deepspeedinference},
smart GPU--CPU swapping~\cite{huang2020swapadvisor}, and KV-cache management
for long-context batching~\cite{kwon2023pagedattention}.
These techniques deliver substantial gains in some configurations yet show
diminishing returns---or even performance inversion---in others.
This variability raises a fundamental question that incremental optimization
studies do not answer: \emph{Is hierarchical memory orchestration a continuously
optimizable engineering problem, or does it admit intrinsic boundaries beyond
which additional coordination cannot improve---and can worsen---end-to-end
latency?}

\paragraph{This work.}
We answer this question with a regime-based view.
We show that multi-tier inference systems exhibit \emph{phase-like} operational
regimes governed by two dimensionless ratios: the fast-memory residency ratio
$R_C = C_{\mathrm{fast}}/W$ (capacity balance) and the transfer-pressure ratio
$R_B = B_{\mathrm{slow}}/D$ (bandwidth balance).
Crucially, transitions between regimes are \emph{discontinuous}: modest changes
in model size, batch size, or available bandwidth abruptly shift the dominant
contributor to latency from locality effects, to coordination overhead, to
physical I/O limits.
Orchestration strategies that appear universally beneficial within one regime
therefore fail predictably---and for principled reasons---outside analytically
identifiable boundaries.

The analogy to physical phase transitions is deliberate.
Just as a material crosses a phase boundary when a control parameter (temperature,
pressure) passes a critical threshold, a hierarchical memory system crosses a
regime boundary when $R_C$ or $R_B$ falls below a critical value $\theta_C$ or
$\theta_B$. Below these thresholds, the dominant latency term changes
discontinuously, and no amount of coordination can compensate for the violated
physical constraint.

\paragraph{Scope and role of \textsc{Orion}.}
\textsc{Orion} is introduced as an \emph{experimental probe} to expose and
measure regime-dependent behavior, not as a new state-of-the-art inference
engine.
Our goal is to test the widespread assumption that deeper orchestration is always
beneficial.
We show that even careful coordination cannot prevent systems from crossing
intrinsic boundaries where additional orchestration becomes additive overhead
rather than a remedy.
Improvements observed with \textsc{Orion} should be interpreted as evidence that
the operating point lies within a coordination-dominated regime; the predictable
failure modes beyond regime boundaries illustrate limits that no orchestration
policy can eliminate.

\paragraph{Generality.}
Although we use large language model serving as a canonical stress test, the
regime structure is not LLM-specific.
Whenever an active working set cannot fully reside in fast memory---as in vision
and multimodal inference, retrieval-augmented generation (RAG), data-intensive
analytics, and HPC-style scientific workflows---performance is governed by the
same intrinsic balance between residency, sustained transfer bandwidth, and
synchronization overhead.
The regime framework therefore provides a general lens for reasoning about
scalability limits in any hierarchical-memory computing system.

\paragraph{Contributions.}
We make three contributions:
\begin{itemize}
  \item \textbf{Regime-dependent principle.}
    We identify three intrinsic operational regimes---capacity-limited,
    coordination-dominated, and I/O-limited---and demonstrate that transitions
    between them are abrupt rather than gradual, analogous to phase transitions
    in physical systems.
  \item \textbf{Minimal analytical model.}
    We derive a structural lower bound on end-to-end inference latency from a
    minimal multi-tier transfer model, yielding analytically interpretable regime
    boundaries expressed as closed-form conditions on $R_C$ and $R_B$.
  \item \textbf{Multi-platform experimental validation.}
    Using \textsc{Orion} as a measurement instrument, we validate predicted
    transitions across language and vision workloads on GPU (NVIDIA A100),
    TPU~v4, Inferentia2, AMD MI250, and Intel Optane-PMem systems, confirming
    that regime boundaries are hardware-agnostic within $\pm 0.05$ of predicted
    thresholds.
\end{itemize}

Together, these results reframe hierarchical memory coordination as a problem
of \emph{regime-aware design} rather than continuous tuning, providing a
principled framework for understanding scalability, stability, and energy
efficiency in memory-bound AI inference.

\FloatBarrier
